% Draft content placeholder only; final paper-facing artwork is intentionally absent.
\begin{figure*}[t]
\centering
\fbox{\parbox[c][31mm][c]{0.96\textwidth}{\small
\textbf{Draft Fig. 2 specification.} Left: absolute Object-126 success curves
for stale arm/fresh gripper and fresh arm/stale gripper at 0.10, 0.20, 0.40,
0.60, 0.80, 1.00, and 1.60~s. Mark 1.00~s as the originally frozen lag, not
an optimum. Right: five block-matched 1.00-s values, with distinct encodings
for historical anchors and later probes: translation stale 8.73\%, full arm
stale 9.52\%, rotation stale 42.06\%, Fresh 44.44\%, and gripper stale
64.29\%. Annotate translation-minus-rotation as $-33.33$ points.}}
\caption{\textbf{Age dependence and within-arm component structure on the
Object development cohort.} The fresh-arm/stale-gripper branch remains above
the stale-arm/fresh-gripper branch at every tested age, with observed
separations from $+21.43$ to $+54.76$ points; both curves are non-monotone.
The 1.00-s separation is the largest observed on the grid, but no
lag-versus-lag inference establishes a distinct peak. Historical anchors and
reviewer-directed probes are exactly block matched but were generated at
different study stages. The right panel is therefore a comparative
characterization, not a single five-condition preregistered factorial. The
stale-gripper effect magnitude should not be generalized across cohorts.}
\label{fig:development}
\end{figure*}
